\documentclass[12pt]{article}

\usepackage{sbc-template}
\usepackage{graphicx,url}
\usepackage[utf8]{inputenc}
\usepackage[brazil]{babel}
\usepackage[latin1]{inputenc}  
\usepackage{booktabs}

\usepackage{float}
\usepackage{tikz}
\usetikzlibrary{positioning,arrows.meta,fit,backgrounds,calc}

     
\sloppy

\title{Evaluation of the Digital Accessibility of Institutional Portals and Electronic Ombudsman Channels in the Municipalities of Mato Grosso}

\author{José Augusto Pereira\inst{1}, Giuseph de Luka Giangareli\inst{1}, Claudemir Marcelino Avelino\inst{1}, \\ Lucas Rocha Bariani\inst{1} Jean Zahn\inst{1} Wellynton Diniz \inst{1}}

\address{Department of Computer Engineering -- Faculdade de Tecnologia de Sinop (FASTECH)\\
    Estrada Claudete 442a -- 91.501-970 -- Sinop -- MT -- Brazil
  \email{\{jpereira, ggiangareli, cavelino, lbariani, jzahn, wdiniz\}@enc.fastech.edu.br}
}

\begin{document} 

\maketitle

\begin{abstract}
\textbf{Research Context:} Digital accessibility is a legal and social requirement for Brazilian government portals, particularly for ombudsman (ouvidoria) services, an essential channel for interaction between citizens and municipal public administration.
\textbf{Scientific and/or Practical Problem:} There is no systematic evidence on the level of digital accessibility of the institutional portals of Mato Grosso's municipalities, nor on whether this standard is maintained in their electronic ombudsman channels, which are often hosted on separate domains or systems.
\textbf{Proposed Solution and/or Analysis:} This study proposes evaluating, using the automated tool AMAWeb (based on WCAG 2.1), the official portals and their corresponding ombudsman channels for all municipalities in the state, comparing the scores obtained to identify consistency or disparity between the two environments.
\textbf{Related IS Theory:} The study draws on the DeLone \& McLean IS Success Model, in which system quality, here operationalized through compliance with accessibility guidelines, is a relevant antecedent of effective use and citizen satisfaction with public digital services.
\textbf{Research Method:} A descriptive and comparative study, with a quantitative approach, involving a census-based automated evaluation of Mato Grosso's municipalities, classification of ombudsman channels by type, data recorded in a standardized spreadsheet (Municipality, Population, URLs, Scores, Errors, date/time), and evidence collection through screenshots and AMAWeb reports.
\textbf{Summary of Results:} The study is expected to identify the percentage of municipalities with accessible portals and ombudsman channels, map the most common types of ombudsman solutions, and verify whether a correlation exists between population size and level of accessibility.
\textbf{Contributions and Impact to IS area:} This work offers an unprecedented diagnosis of public digital accessibility in Mato Grosso, supports public managers in prioritizing improvements, and expands the empirical application of AMAWeb as an evaluation instrument in Information Systems research for the public sector.
\end{abstract}
     
\begin{resumo} 
\textbf{Research Context:} A acessibilidade digital é um requisito legal e social para portais governamentais brasileiros, especialmente em serviços de ouvidoria, canal essencial de interação entre cidadão e administração pública municipal.
\textbf{Scientific and/or Practical Problem:} Não há evidências sistemáticas sobre o nível de acessibilidade dos portais institucionais das prefeituras de Mato Grosso nem sobre se esse padrão se mantém nos canais eletrônicos de ouvidoria, muitas vezes hospedados em domínios ou sistemas distintos.
\textbf{Proposed Solution and/or Analysis:} Propõe-se avaliar, com a ferramenta automatizada AMAWeb (baseada nas WCAG 2.1), os portais oficiais e os respectivos canais de ouvidoria de todos os municípios do estado, comparando as notas obtidas para identificar consistência ou disparidade entre os dois ambientes.
\textbf{Related IS Theory:} O estudo se apoia no DeLone & McLean IS Success Model, no qual a qualidade do sistema, aqui operacionalizada pela conformidade com diretrizes de acessibilidade, é um antecedente relevante do uso efetivo e da satisfação dos cidadãos com serviços digitais públicos.
\textbf{Research Method:} Pesquisa descritiva e comparativa, de abordagem quantitativa, com avaliação automatizada censitária dos municípios de Mato Grosso, classificação dos canais de ouvidoria por tipo, registro em planilha padronizada (Município, População, URLs, Notas, Erros, data/hora) e evidenciação por captura de tela e relatório do AMAWeb.
\textbf{Summary of Results:} Espera-se identificar o percentual de municípios com portais e ouvidorias acessíveis, mapear os tipos de solução de ouvidoria mais comuns e verificar se existe correlação entre porte populacional e nível de acessibilidade.
\textbf{Contributions and Impact to IS area:} O trabalho oferece um diagnóstico inédito da acessibilidade digital pública em Mato Grosso, subsidia gestores na priorização de melhorias e amplia a aplicação empírica do AMAWeb como instrumento de avaliação em Sistemas de Informação para o setor público.
\end{resumo}


\section{Introdução}

A acessibilidade digital e um requisito legal e social para portais governamentais brasileiros, desta forma entender se os portais municipais assim como suas ouvidorias atende a esses critérios e importante para garantir a equidade de todos os cidadães e também para garantir a interação entre todos os cidadãos com acesso a internet e administração publica. Em seu trabalho \cite{wcge} realizou uma analise dos portais governamentais apresentou um dado interessante sobre a consistência da pontuação de acessibilidade nas paginas secundaria, desta forma como o trabalho nasce de um grupo que tem como foco governo digital e ouvidorias municipais como objeto de estudo seria interessante registar se há alguma consistência entre o portal da prefeitura, pagina da ouvidoria, e portal da câmara legislativa municipal. 

A escolha do estado de Mato Grosso ocorre por suas características geográficas, demográficas e administrativas, o que o torna relevante para a analise da acessibilidade dos serviços públicos, disponibilizados em meios digitais. O estado possui mais de 900 mil km² e baixa densidade demográfica, tendo seus municípios distribuídos em um  amplo território e diferentes faixas populacionais. Nesse contexto, os canais digitais da administração pública municipal podem assumir papel importante na aproximação entre cidadãos e poder público, possibilitando o acesso a informações, serviços e mecanismos de participação sem a necessidade de deslocamento físico.

Ao analisamos os município mato-grossenses permite entender como a acessibilidade digital e pedimentada, permitindo investigar possíveis diferenças de implantação relacionadas a a característico as distintas tanto populacional quanto socioeconômicas. Essa delimitação favorece a realização de uma avaliação comparativa e padronizada dos portais institucionais das prefeituras e de seus respectivos canais eletrônicos de ouvidoria.

O escopo do trabalho se alinhada às iniciativas atuais de transformação digital da administração pública. No âmbito estadual, Mato Grosso instituiu sua Agenda Estratégica Digital para o período de 2023 a 2027, que apresentava como uma de suas metas a "Adaptar todos os serviços digitais do Poder Executivo estadual ao padrão de acessibilidade e-MAG, até 2026". \cite{matogrosso:2023} contemplando ações relacionadas à modernização, inovação e transformação digital da administração pública. Dessa forma, avaliar a acessibilidade dos canais digitais municipais pode contribuir para identificar barreiras existentes e fornecer um diagnóstico sobre as condições de acesso digital oferecidas aos cidadãos no estado.


\section{Trabalhos relacionados}
A Acessibilidade em portais governamentais tem sido investigada de sob diferentes perspectivas, abrangendo desde órgãos da administração publica federal, estadual e municipal. De modo geral os estudos apontam a existência de dificuldades recorrentes no atendimento  ás recomendações de acessibilidade digital, mesmo diante da existência de normas e modelos específicos para o contexto governamental brasileiro. 

\cite{sbsi} realizaram em 2015 um estudo voltados à aplicação de padrões de acessibilidade em portais governamentais brasileiros. Os autores avaliaram 39 sítios relacionados aos ministérios do Governo Federal, considerando padrões e recomendações como o Modelo de Acessibilidade em Governo Eletrônico (e-MAG) e padrões web estabelecidos pelo W3C. Os resultados evidenciaram a presença de erros capazes de comprometer o acesso às informações e aos serviços digitais, indicando que, apesar da existência de diretrizes governamentais, sua aplicação ainda ocorria de maneira insuficiente.

Posteriormente \cite{10.1145/3411564.3411656} retornam a investigação sobre acessibilidade no Governo Federal, com o objetivo de analisar a evolução dos portais governamentais em relação às avaliações anteriores,  possibilitou observar alterações ocorridas nos portais ao longo do tempo e verificar em que medida as iniciativas adotadas pelo Governo Federal refletiram em melhorias na acessibilidade. Demostrando a importância não apenas de identificar problemas pontuais, mas também de realizar avaliações periódicas para acompanhar a evolução da acessibilidade dos sistemas de informação governamentais.

Essa linha de investigação foi novamente explorada por  \cite{10.1145/3658271.3658276}  realizou um estudo de replicação sobre acessibilidade digital em 36 portais do Governo Federal deses o mesmo identificou que foi possível analisar a evolução de apenas 17 portais, ele Utilizou o ASES, para verificar a conformidade com o eMAG e  a ferramenta Wappalyzer. Destaca no texto a questão do padrão E-mag ser de 2014 quando o WCAG era de 2023 ano do estudo.

No âmbito dos municípios \cite{sbsiAmorim_Menezes} propuseram uma metodologia de avaliação de Portais da Transparências Municipais,com base nas normas da série ISO/IEC 25000 e na legislação brasileira. A pesquisa demonstra a viabilidade da utilização de procedimentos padronizados para analisar um número significativo de portais públicos, reduzindo a subjetividade do processo avaliativo. 

Em um estudo que aborda especificamente as ouvidorias públicas, \cite{wcge} investigam o uso de Tecnologias da Informação e Comunicação nas Ouvidorias Gerais Municipais do Estado de Mato Grosso. A pesquisa contemplou 20 municípios e analisou os canais de comunicação existentes e a disponibilidade de informações aos cidadãos. Os resultados mostraram que apenas 15\% das ouvidorias avaliadas atendiam a todos os requisitos considerados para acesso e disponibilização de informações. Por abordar as ouvidorias municipais, o trabalho fornece uma referência para compreender as limitações existentes na interação digital entre cidadãos e administrações municipais no estado.

Mais recentemente, \cite{eres}, realizaram uma avaliação preliminar dos sites das prefeituras onde estão localizadas os Campi da Unipampa com apoio do Avaliador Avaliador e Simulador de Acessibilidade em Sítios (ASES) para verificar a conformidade das páginas com recomendações do e-MAG. Além de aproximar a discussão do contexto municipal, a pesquisa demonstra a aplicabilidade de ferramentas automáticas para identificar problemas de acessibilidade em portais de prefeituras. 

O que esses trabalhos apresentam é que mesmo tendo leis e normas para garantir a acessibilidade digital, ao realizar avaliações foi identificado erros e falhas as quais não garante a equidade ao acesso da informação. No estudo de forma preliminar busca gerar um Ranking que permita a evolução da analise para entender a evolução do Governo Digital nos municípios que compõem o Estado de Mato Grosso. 
\section{Metodologia}

O estudo parte da identificação dos portais das prefeituras e dos canais de ouvidoria dos 142 municípios do estado. Cada portal e cada ouvidoria são avaliados com duas ferramentas automatizadas baseadas nas WCAG 2.1~\cite{wcag21}, o AMAWeb e o AccessMonitor, e o relatório de cada avaliação é armazenado em JSON e PDF. Para tornar a coleta mais ágil, a submissão das URLs às ferramentas é automatizada por \textit{scripts}. A Figura~\ref{fig:metodologia} resume o processo, organizado em três fases: levantamento, avaliação automatizada e consolidação e análise.

% ============================================================
% Figura: Processo metodológico
% Requer no preâmbulo:
%   \usepackage{tikz}
%   \usetikzlibrary{positioning,arrows.meta,fit,backgrounds,calc}
% Uso no corpo do texto:  % ============================================================
% Figura: Processo metodológico
% Requer no preâmbulo:
%   \usepackage{tikz}
%   \usetikzlibrary{positioning,arrows.meta,fit,backgrounds,calc}
% Uso no corpo do texto:  % ============================================================
% Figura: Processo metodológico
% Requer no preâmbulo:
%   \usepackage{tikz}
%   \usetikzlibrary{positioning,arrows.meta,fit,backgrounds,calc}
% Uso no corpo do texto:  \input{figura-metodologia}
% ============================================================
\begin{figure}[t]
\centering
\resizebox{\textwidth}{!}{%
\begin{tikzpicture}[
  node distance=5mm and 12mm,
  etapa/.style={rectangle, rounded corners=2pt, draw=black!70, fill=white,
                text width=5.2cm, align=center, font=\small,
                minimum height=1.2cm, inner sep=4pt},
  apoio/.style={etapa, dashed, fill=black!5},
  seta/.style={-{Stealth[length=2.5mm]}, thick, draw=black!70},
  fase/.style={draw=black!40, fill=black!4, rounded corners=4pt, inner sep=3mm},
  titulo/.style={font=\bfseries\small, minimum height=9mm, anchor=north, text width=5.2cm, align=center}
]

% ---------------- Fase 1: Levantamento ----------------
\node[titulo] (t1) {Fase 1\\Levantamento};
\node[etapa, below=of t1] (e1)
  {Identificação dos portais de prefeituras e câmaras\\{\footnotesize(Radar da Transparência; 314 URLs)}};
\node[etapa, below=of e1] (e2)
  {Localização manual dos canais de ouvidoria\\{\footnotesize(142 municípios)}};
\node[etapa, below=of e2] (e3)
  {Classificação do tipo de ouvidoria\\[2pt]
   {\scriptsize integrada ao portal $\cdot$ sistema próprio (domínio/subdomínio)
    $\cdot$ terceirizada $\cdot$ Fala.BR
    $\cdot$ não localizada $\cdot$ indisponível}};
    
% ---------------- Fase 2: Avaliação ----------------
\node[titulo, right=of t1] (t2) {Fase 2\\Avaliação automatizada};
\node[apoio, below=of t2] (e4)
  {\textit{Scripts} em Python\\{\footnotesize(requisições diretas ao serviço de avaliação)}};
\node[etapa, below=of e4] (e5)
  {Avaliação com AMAWeb e AccessMonitor\\{\footnotesize (WCAG 2.1) prefeituras, câmaras e ouvidorias; até cinco tentativas}};
\node[apoio, below=of e5] (e5b)
  {Contingência: envio do HTML baixado\\{\footnotesize(páginas bloqueadas por serviços de proteção)}};
\node[etapa, below=of e5b] (e6)
  {Registro de evidências e validação\\{\footnotesize JSON bruto, PDF e \textit{hash}; uma avaliação por município, ambiente e ferramenta}};

% ---------------- Fase 3: Consolidação e análise ----------------
\node[titulo, right=of t2] (t3) {Fase 3\\Consolidação e análise};
\node[etapa, below=of t3] (e7)
  {Planilha consolidada\\{\footnotesize URL, ferramenta, situação, nota, falhas, data/hora e \textit{hash}}};
\node[apoio, below=of e7] (e8)
  {Dados complementares\\{\footnotesize População (IBGE, 2026), PIB \textit{per capita} (2023), IDHM (2010)}};
\node[etapa, below=of e8] (e9)
  {Estatística descritiva (QP1)\\{\footnotesize média, mediana, dispersão e faixas de nota}};
\node[etapa, below=of e9] (e10)
  {Comparação entre ambientes e por tipo de ouvidoria (QP2 e QP3)\\{\footnotesize diferenças prefeitura $-$ ouvidoria e prefeitura $-$ câmara; Wilcoxon pareado}};
\node[etapa, below=of e10] (e11)
  {Correlação com indicadores (QP4), erros e ranqueamento\\{\footnotesize Spearman; erros mais frequentes; média das três notas}};

% ---------------- Caixas das fases (mesma altura) ----------------
\begin{scope}[on background layer]
  \node[fase, fit=(t1)(e1)(e3)(e3|-e11.south)] (f1) {};
  \node[fase, fit=(t2)(e4)(e6)(e6|-e11.south)] (f2) {};
  \node[fase, fit=(t3)(e7)(e11)] (f3) {};
\end{scope}

% ---------------- Setas ----------------
\draw[seta] (e1) -- (e2);
\draw[seta] (e2) -- (e3);
\draw[seta] (e4) -- (e5);
\draw[seta] (e5) -- (e5b);
\draw[seta] (e5b) -- (e6);
\draw[seta] (e7) -- (e8);
\draw[seta] (e8) -- (e9);
\draw[seta] (e9) -- (e10);
\draw[seta] (e10) -- (e11);

\draw[seta] (f1.east) -- node[above, font=\scriptsize, align=center]{URLs}   (f2.west);
\draw[seta] (f2.east) -- node[above, font=\scriptsize, align=center]{notas}  (f3.west);

\end{tikzpicture}%
}
\caption{Processo metodológico adotado na avaliação da acessibilidade dos portais das prefeituras, dos portais das câmaras e dos canais de ouvidoria dos municípios de Mato Grosso.}
\label{fig:metodologia}
\end{figure}
% ============================================================
\begin{figure}[t]
\centering
\resizebox{\textwidth}{!}{%
\begin{tikzpicture}[
  node distance=5mm and 12mm,
  etapa/.style={rectangle, rounded corners=2pt, draw=black!70, fill=white,
                text width=5.2cm, align=center, font=\small,
                minimum height=1.2cm, inner sep=4pt},
  apoio/.style={etapa, dashed, fill=black!5},
  seta/.style={-{Stealth[length=2.5mm]}, thick, draw=black!70},
  fase/.style={draw=black!40, fill=black!4, rounded corners=4pt, inner sep=3mm},
  titulo/.style={font=\bfseries\small, minimum height=9mm, anchor=north, text width=5.2cm, align=center}
]

% ---------------- Fase 1: Levantamento ----------------
\node[titulo] (t1) {Fase 1\\Levantamento};
\node[etapa, below=of t1] (e1)
  {Identificação dos portais de prefeituras e câmaras\\{\footnotesize(Radar da Transparência; 314 URLs)}};
\node[etapa, below=of e1] (e2)
  {Localização manual dos canais de ouvidoria\\{\footnotesize(142 municípios)}};
\node[etapa, below=of e2] (e3)
  {Classificação do tipo de ouvidoria\\[2pt]
   {\scriptsize integrada ao portal $\cdot$ sistema próprio (domínio/subdomínio)
    $\cdot$ terceirizada $\cdot$ Fala.BR
    $\cdot$ não localizada $\cdot$ indisponível}};
    
% ---------------- Fase 2: Avaliação ----------------
\node[titulo, right=of t1] (t2) {Fase 2\\Avaliação automatizada};
\node[apoio, below=of t2] (e4)
  {\textit{Scripts} em Python\\{\footnotesize(requisições diretas ao serviço de avaliação)}};
\node[etapa, below=of e4] (e5)
  {Avaliação com AMAWeb e AccessMonitor\\{\footnotesize (WCAG 2.1) prefeituras, câmaras e ouvidorias; até cinco tentativas}};
\node[apoio, below=of e5] (e5b)
  {Contingência: envio do HTML baixado\\{\footnotesize(páginas bloqueadas por serviços de proteção)}};
\node[etapa, below=of e5b] (e6)
  {Registro de evidências e validação\\{\footnotesize JSON bruto, PDF e \textit{hash}; uma avaliação por município, ambiente e ferramenta}};

% ---------------- Fase 3: Consolidação e análise ----------------
\node[titulo, right=of t2] (t3) {Fase 3\\Consolidação e análise};
\node[etapa, below=of t3] (e7)
  {Planilha consolidada\\{\footnotesize URL, ferramenta, situação, nota, falhas, data/hora e \textit{hash}}};
\node[apoio, below=of e7] (e8)
  {Dados complementares\\{\footnotesize População (IBGE, 2026), PIB \textit{per capita} (2023), IDHM (2010)}};
\node[etapa, below=of e8] (e9)
  {Estatística descritiva (QP1)\\{\footnotesize média, mediana, dispersão e faixas de nota}};
\node[etapa, below=of e9] (e10)
  {Comparação entre ambientes e por tipo de ouvidoria (QP2 e QP3)\\{\footnotesize diferenças prefeitura $-$ ouvidoria e prefeitura $-$ câmara; Wilcoxon pareado}};
\node[etapa, below=of e10] (e11)
  {Correlação com indicadores (QP4), erros e ranqueamento\\{\footnotesize Spearman; erros mais frequentes; média das três notas}};

% ---------------- Caixas das fases (mesma altura) ----------------
\begin{scope}[on background layer]
  \node[fase, fit=(t1)(e1)(e3)(e3|-e11.south)] (f1) {};
  \node[fase, fit=(t2)(e4)(e6)(e6|-e11.south)] (f2) {};
  \node[fase, fit=(t3)(e7)(e11)] (f3) {};
\end{scope}

% ---------------- Setas ----------------
\draw[seta] (e1) -- (e2);
\draw[seta] (e2) -- (e3);
\draw[seta] (e4) -- (e5);
\draw[seta] (e5) -- (e5b);
\draw[seta] (e5b) -- (e6);
\draw[seta] (e7) -- (e8);
\draw[seta] (e8) -- (e9);
\draw[seta] (e9) -- (e10);
\draw[seta] (e10) -- (e11);

\draw[seta] (f1.east) -- node[above, font=\scriptsize, align=center]{URLs}   (f2.west);
\draw[seta] (f2.east) -- node[above, font=\scriptsize, align=center]{notas}  (f3.west);

\end{tikzpicture}%
}
\caption{Processo metodológico adotado na avaliação da acessibilidade dos portais das prefeituras, dos portais das câmaras e dos canais de ouvidoria dos municípios de Mato Grosso.}
\label{fig:metodologia}
\end{figure}
% ============================================================
\begin{figure}[t]
\centering
\resizebox{\textwidth}{!}{%
\begin{tikzpicture}[
  node distance=5mm and 12mm,
  etapa/.style={rectangle, rounded corners=2pt, draw=black!70, fill=white,
                text width=5.2cm, align=center, font=\small,
                minimum height=1.2cm, inner sep=4pt},
  apoio/.style={etapa, dashed, fill=black!5},
  seta/.style={-{Stealth[length=2.5mm]}, thick, draw=black!70},
  fase/.style={draw=black!40, fill=black!4, rounded corners=4pt, inner sep=3mm},
  titulo/.style={font=\bfseries\small, minimum height=9mm, anchor=north, text width=5.2cm, align=center}
]

% ---------------- Fase 1: Levantamento ----------------
\node[titulo] (t1) {Fase 1\\Levantamento};
\node[etapa, below=of t1] (e1)
  {Identificação dos portais de prefeituras e câmaras\\{\footnotesize(Radar da Transparência; 314 URLs)}};
\node[etapa, below=of e1] (e2)
  {Localização manual dos canais de ouvidoria\\{\footnotesize(142 municípios)}};
\node[etapa, below=of e2] (e3)
  {Classificação do tipo de ouvidoria\\[2pt]
   {\scriptsize integrada ao portal $\cdot$ sistema próprio (domínio/subdomínio)
    $\cdot$ terceirizada $\cdot$ Fala.BR
    $\cdot$ não localizada $\cdot$ indisponível}};
    
% ---------------- Fase 2: Avaliação ----------------
\node[titulo, right=of t1] (t2) {Fase 2\\Avaliação automatizada};
\node[apoio, below=of t2] (e4)
  {\textit{Scripts} em Python\\{\footnotesize(requisições diretas ao serviço de avaliação)}};
\node[etapa, below=of e4] (e5)
  {Avaliação com AMAWeb e AccessMonitor\\{\footnotesize (WCAG 2.1) prefeituras, câmaras e ouvidorias; até cinco tentativas}};
\node[apoio, below=of e5] (e5b)
  {Contingência: envio do HTML baixado\\{\footnotesize(páginas bloqueadas por serviços de proteção)}};
\node[etapa, below=of e5b] (e6)
  {Registro de evidências e validação\\{\footnotesize JSON bruto, PDF e \textit{hash}; uma avaliação por município, ambiente e ferramenta}};

% ---------------- Fase 3: Consolidação e análise ----------------
\node[titulo, right=of t2] (t3) {Fase 3\\Consolidação e análise};
\node[etapa, below=of t3] (e7)
  {Planilha consolidada\\{\footnotesize URL, ferramenta, situação, nota, falhas, data/hora e \textit{hash}}};
\node[apoio, below=of e7] (e8)
  {Dados complementares\\{\footnotesize População (IBGE, 2026), PIB \textit{per capita} (2023), IDHM (2010)}};
\node[etapa, below=of e8] (e9)
  {Estatística descritiva (QP1)\\{\footnotesize média, mediana, dispersão e faixas de nota}};
\node[etapa, below=of e9] (e10)
  {Comparação entre ambientes e por tipo de ouvidoria (QP2 e QP3)\\{\footnotesize diferenças prefeitura $-$ ouvidoria e prefeitura $-$ câmara; Wilcoxon pareado}};
\node[etapa, below=of e10] (e11)
  {Correlação com indicadores (QP4), erros e ranqueamento\\{\footnotesize Spearman; erros mais frequentes; média das três notas}};

% ---------------- Caixas das fases (mesma altura) ----------------
\begin{scope}[on background layer]
  \node[fase, fit=(t1)(e1)(e3)(e3|-e11.south)] (f1) {};
  \node[fase, fit=(t2)(e4)(e6)(e6|-e11.south)] (f2) {};
  \node[fase, fit=(t3)(e7)(e11)] (f3) {};
\end{scope}

% ---------------- Setas ----------------
\draw[seta] (e1) -- (e2);
\draw[seta] (e2) -- (e3);
\draw[seta] (e4) -- (e5);
\draw[seta] (e5) -- (e5b);
\draw[seta] (e5b) -- (e6);
\draw[seta] (e7) -- (e8);
\draw[seta] (e8) -- (e9);
\draw[seta] (e9) -- (e10);
\draw[seta] (e10) -- (e11);

\draw[seta] (f1.east) -- node[above, font=\scriptsize, align=center]{URLs}   (f2.west);
\draw[seta] (f2.east) -- node[above, font=\scriptsize, align=center]{notas}  (f3.west);

\end{tikzpicture}%
}
\caption{Processo metodológico adotado na avaliação da acessibilidade dos portais das prefeituras, dos portais das câmaras e dos canais de ouvidoria dos municípios de Mato Grosso.}
\label{fig:metodologia}
\end{figure}

\subsection{População e fonte dos endereços}

A população do estudo compreende os 142 municípios de Mato Grosso. Para cada município, são consideradas duas unidades de avaliação: o portal institucional da prefeitura e o canal eletrônico de ouvidoria.

Os endereços dos portais são obtidos no Radar da Transparência Pública~\cite{radar}, painel do Programa Nacional de Transparência Pública que reúne os entes públicos avaliados e seus endereços eletrônicos. São extraídos os entes vinculados aos municípios mato-grossenses no ciclo de 2025. Cada endereço é classificado como prefeitura, câmara ou outro entre, com base no domínio (por exemplo, \texttt{.leg.br} para câmaras) e no título da página. As análises deste artigo restringem-se aos portais das prefeituras.

\subsection{Identificação e classificação das ouvidorias}

Os canais de ouvidoria são localizados manualmente, por navegação a partir do portal oficial de cada município. Para cada município, registram-se em planilha:
\begin{itemize}
  \item o portal oficial e a URL de partida;
  \item a URL final alcançada, ou a última tentada;
  \item o título da página;
  \item o caminho de navegação percorrido (por exemplo, \textit{Portal $\rightarrow$ Ouvidoria $\rightarrow$ Novo protocolo});
  \item uma descrição da evidência encontrada;
  \item a situação do canal: formulário confirmado, canal sem formulário, portal intermediário, página sem conteúdo visível, erro de acesso ou não verificável.
\end{itemize}
A URL submetida às ferramentas é a URL final, isto é, a página na qual o cidadão registra a manifestação.

Cada ouvidoria é classificada em um dos seguintes tipos, com base no endereço de hospedagem e na situação registrada:
(i)~integrada ao portal, quando está no mesmo domínio do portal da prefeitura;
(ii)~sistema próprio, quando está em domínio ou subdomínio distinto da prefeitura;
(iii)~terceirizada, quando está hospedada em domínio de fornecedor;
(iv)~Fala.BR, plataforma da Controladoria-Geral da União;
(v)~não localizada; e
(vi)~indisponível no momento da coleta.
% O tipo "formulário simples" previsto originalmente não é distinguível pelo endereço;
% exige classificação manual. Decidir se entra ou não.

\subsection{Instrumentos de avaliação}

São utilizadas duas ferramentas automatizadas de avaliação de acessibilidade baseadas nas WCAG 2.1.

O AMAWeb~\cite{amaweb}, desenvolvido pela Universidade Federal de São Paulo (UNIFESP) e pelo Instituto Federal do Rio Grande do Sul (IFRS), avalia uma página a partir de sua URL ou de seu código HTML. Ele atribui uma nota de 1 a 10 e classifica cada prática verificada como aceitável, não aceitável ou a verificar manualmente.

O AccessMonitor~\cite{accessmonitor}, validador do Observatório Português da Acessibilidade Web, mantido pela Agência para a Modernização Administrativa (AMA) de Portugal, executa o motor QualWeb~\cite{qualweb}. O QualWeb aplica regras ACT, técnicas WCAG 2.1 e boas práticas, e sobre esses resultados a ferramenta calcula uma nota de 1 a 10. Por ter código aberto, a ferramenta é executada em instância local, instalada a partir dos repositórios oficiais, nas seguintes versões: serviço \texttt{accessmonitor-docker} (junho de 2026), QualWeb \textit{core}~0.8.11, \textit{act-rules}~0.8.0, \textit{wcag-techniques}~0.4.7, \textit{best-practices}~0.7.12 e conjunto de regras \texttt{accessmonitor-rulesets}~1.0.4.

As duas ferramentas devolvem relatórios com a mesma estrutura: nota, práticas classificadas pelos três vereditos, contagem de elementos HTML e título da página. As notas de cada ferramenta são analisadas separadamente.

\subsection{Procedimento de coleta}

A coleta é automatizada por \textit{scripts} em Python, um para cada ferramenta:
\begin{itemize}
  \item \textbf{AMAWeb:} o \textit{script} consulta o serviço HTTP utilizado pela própria interface web do avaliador, que recebe a URL codificada e devolve o relatório em JSON. Esse serviço não possui documentação pública formal.
  \item \textbf{AccessMonitor:} o \textit{script} consulta a instância local, que recebe a URL codificada em Base64 ou o código HTML da página. As requisições informam a origem autorizada no cabeçalho \texttt{Referer}, conforme exigido pela instância.
\end{itemize}

Para cada URL, os \textit{scripts} preservam o JSON devolvido pela ferramenta, sem alterações, e geram um relatório em PDF. As requisições são executadas com três avaliações simultâneas e limite de tempo por avaliação. Cada execução grava seus resultados em um diretório identificado pela data e pela hora de início (\texttt{AAAA-MM-DD/HH-MM-SS}), acompanhado de uma planilha CSV com o resultado de cada URL e de um manifesto com os parâmetros utilizados. As URLs sem resultado válido são reenviadas em rodadas sucessivas, até o limite de cinco tentativas.

As páginas que bloqueiam o acesso direto da ferramenta são avaliadas por um modo alternativo: o HTML da página é baixado com cabeçalhos de um navegador comum e submetido ao modo de avaliação por código HTML. Nesse modo, os \textit{scripts} da página não são executados, por isso o modo de envio (URL ou HTML) é registrado em cada avaliação. Sites com certificado TLS inválido não são avaliados, pois a validação do certificado não é contornada.

\subsection{Critérios de validade}

Uma avaliação é considerada válida quando o documento analisado corresponde à página solicitada. São descartadas as avaliações de:
(i)~páginas de desafio ou de acesso negado de serviços de proteção, identificadas por marcadores no título (``\textit{Just a moment}'', ``\textit{Attention Required!}'', ``\textit{Access denied}'');
(ii)~páginas de erro HTTP, com código igual ou superior a 400; e
(iii)~documentos com menos de dez elementos HTML.
As ouvidorias classificadas como não localizadas ou indisponíveis não integram a análise das notas, pois a URL correspondente não é um canal funcional. Quando uma mesma URL tem mais de uma avaliação válida na mesma ferramenta, mantém-se a feita pela URL em detrimento da feita pelo HTML e, entre essas, a mais recente.

\subsection{Análise dos dados}

Os relatórios são consolidados em uma planilha com cinco abas: câmaras, prefeituras, ouvidorias, análise e erros. As três primeiras registram URL, ferramenta, situação da página, nota informada, nota considerada na média, data e hora, quantidade de práticas com falha, aviso e sucesso, e o \textit{hash} SHA-256 do arquivo de origem, que permite verificar a correspondência com o relatório bruto. A aba de análise reúne os ranqueamentos por município, e a de erros relaciona as falhas por URL e ferramenta.

A unidade de análise é o município. Para cada ferramenta, cada município tem duas notas: a do portal da prefeitura e a da ouvidoria. Os indicadores são a média simples entre as duas notas e a diferença entre a nota do portal e a da ouvidoria. A diferença reflete o percurso do cidadão, que normalmente chega à ouvidoria a partir do portal.

As notas são descritas por média, mediana, desvio-padrão, mínimo e máximo, para cada ambiente e ferramenta. Os municípios são ranqueados pela média simples entre as notas do portal e da ouvidoria, e são apresentados os dez mais bem e os dez mais mal posicionados. As notas são ainda relacionadas aos indicadores municipais, por meio do coeficiente de correlação de Spearman e de médias por faixa populacional: população estimada pelo IBGE (2026), PIB \textit{per capita} e IDHM (2010).

\section{Resultados e discussões}

\subsection{Caracterização dos dados}

A coleta foi realizada entre 20 e 22 de setembro de 2026. Do Radar da Transparência foram extraídas 314 URLs únicas, entre portais de prefeituras, de câmaras municipais e de outros entes, como consórcios intermunicipais. A localização das ouvidorias foi feita em 21 de setembro de 2026, e o formulário de manifestação foi confirmado em 123 dos 142 municípios (86,6\%). Nos demais, encontramos canais oficiais sem formulário (3), portais intermediários que exigiam seleção de serviço, autenticação ou solicitação prévia (8), páginas sem conteúdo visível (2), erros de acesso, como HTTP 404 e 504, falha de DNS, certificado inválido e bloqueio (5), e um caso não verificável. Como alguns municípios compartilham a mesma plataforma, os 142 registros correspondem a 137 URLs únicas. Por exemplo, Araguainha, Barra do Garças, Cláudia, Paranaíta, Sorriso e Várzea Grande apontam para o mesmo endereço de registro no Fala.BR: \url{https://falabr.cgu.gov.br/v2/manifestacoes/registrar}. São seis registros municipais, mas não seis páginas distintas.

A Tabela~\ref{tab:amostra} resume a cobertura da coleta por município e por ferramenta. Considerando as duas ferramentas em conjunto, portal e ouvidoria foram avaliados em 115 municípios (81,0\%). Em 11 municípios (7,7\%) a ouvidoria não foi localizada ou estava indisponível no momento da coleta, e em 16 um dos ambientes, ou ambos, ficou sem avaliação válida. Por ferramenta, o AccessMonitor avaliou os dois ambientes em 114 municípios e o AMAWeb em 98, sobretudo por causa das falhas persistentes do AMAWeb: 47 dos 314 endereços permaneceram sem resposta após cinco tentativas. Essas falhas não decorreram de limitação de requisições: no mesmo período, URLs já avaliadas foram reenviadas e processadas normalmente, enquanto os endereços com falha repetiram o erro também na interface web do avaliador, inclusive quando submetidos pelo código HTML. Trata-se, portanto, de um erro interno no processamento dessas páginas, o que motivou a inclusão do AccessMonitor como segunda ferramenta.

\begin{table}[ht]
\centering
\caption{Cobertura da coleta nos municípios de Mato Grosso}
\label{tab:amostra}
\begin{tabular}{lrr}
\toprule
\textbf{Situação} & \textbf{Municípios} & \textbf{\%} \\
\midrule
Portal e ouvidoria avaliados            & 115 & 81,0 \\
Apenas portal avaliado                  &   1 &  0,7 \\
Apenas ouvidoria avaliada               &   8 &  5,6 \\
Ouvidoria não localizada                &   3 &  2,1 \\
Ouvidoria indisponível no momento       &   8 &  5,6 \\
Portal e ouvidoria sem avaliação válida &   7 &  4,9 \\
\midrule
\textbf{Total}                          & \textbf{142} & \textbf{100} \\
\bottomrule
\multicolumn{3}{l}{\footnotesize Ambiente avaliado: avaliação válida em ao menos uma das ferramentas.}
\end{tabular}
\end{table}

Os bloqueios por serviço de proteção foram a principal causa de perda de dados. No AccessMonitor, 51 URLs únicas devolveram página de bloqueio do Cloudflare (33 portais e 18 ouvidorias). Registramos ainda 33 erros internos do serviço ao processar o relatório do QualWeb, 18 esgotamentos de tempo no download do HTML, 5 falhas de DNS, parte delas superada ao repetir a coleta a partir de outra rede, e 3 certificados TLS inválidos. A contingência pelo HTML baixado foi decisiva para a cobertura do AccessMonitor, com 128 avaliações válidas obtidas nesse modo. No AMAWeb, o envio por HTML também permitiu avaliar parte dos portais bloqueados. Conforme o critério de validade definido na metodologia, 15 respostas com menos de dez elementos HTML foram mantidas para auditoria, sem compor as médias ou os ranqueamentos.

As médias obtidas nos dois modos de envio foram próximas: no AMAWeb, 7,73 pela URL e 7,64 pelo HTML; no AccessMonitor, 7,15 e 6,89, respectivamente. Por isso, as avaliações dos dois modos foram analisadas em conjunto.

A Tabela~\ref{tab:tipos} apresenta a frequência de cada tipo de solução de ouvidoria e as notas médias correspondentes.

\begin{table}[ht]
\centering
\caption{Tipos de solução de ouvidoria, notas médias da ouvidoria e diferença média portal $-$ ouvidoria, por ferramenta}
\label{tab:tipos}
\small
\setlength{\tabcolsep}{4pt}
\begin{tabular}{lrrrrrr}
\toprule
 & & & \multicolumn{2}{c}{\textbf{Nota ouvidoria}} & \multicolumn{2}{c}{\textbf{Diferença média}} \\
\cmidrule(lr){4-5}\cmidrule(lr){6-7}
\textbf{Tipo} & \textbf{$n$} & \textbf{\%} & \textbf{AMAWeb} & \textbf{AccessM.} & \textbf{AMAWeb} & \textbf{AccessM.} \\
\midrule
Integrada ao portal                  & 34 & 23,9 & 7,66 & 7,13 & $-$0,09 &    0,07 \\
Sistema próprio                      & 59 & 41,5 & 7,83 & 7,24 &    0,00 &    0,21 \\
Terceirizada                         & 31 & 21,8 & 8,67 & 8,27 & $-$1,63 & $-$2,50 \\
Fala.BR                              &  7 &  4,9 & 8,00 & 8,30 & $-$0,47 & $-$1,30 \\
Não localizada                       &  3 &  2,1 &  --- &  --- &     --- &     --- \\
Indisponível                         &  8 &  5,6 &  --- &  --- &     --- &     --- \\
\midrule
\textbf{Total}                       & \textbf{142} & \textbf{100} & & & & \\
\bottomrule
\multicolumn{7}{l}{\footnotesize Médias calculadas sobre as ouvidorias com avaliação válida em cada ferramenta.} \\
\multicolumn{7}{l}{\footnotesize Diferença negativa indica ouvidoria com nota maior que o portal.}
\end{tabular}
\end{table}

Os sistemas próprios em domínio ou subdomínio distinto são o tipo mais frequente (41,5\%). Chama a atenção que 55 dos 59 sigam o mesmo padrão de endereço (\texttt{ouvidoria.<município>.mt.gov.br/Manifestacao/}). Isso sugere que parte dos canais classificados como próprios é, na verdade, uma mesma solução replicada entre municípios, possivelmente fornecida por uma única empresa e hospedada sob o domínio da prefeitura. As soluções terceirizadas em domínio de fornecedor respondem por 21,8\% e as integradas ao portal por 23,9\%. Apenas 7 municípios (4,9\%) direcionam o cidadão diretamente ao Fala.BR.

\subsection{Acessibilidade dos portais institucionais e canais de ouvidoria}

A Tabela~\ref{tab:descritiva} apresenta a estatística descritiva das notas nos dois ambientes, por ferramenta. Considerou-se acessível o ambiente com nota igual ou superior a XX. %justificar o ponto de corte

\begin{table}[H]
\centering
\caption{Estatística descritiva das notas de acessibilidade, por ambiente e ferramenta}
\label{tab:descritiva}
\small
\begin{tabular}{lrrrr}
\toprule
 & \multicolumn{2}{c}{\textbf{Portal}} & \multicolumn{2}{c}{\textbf{Ouvidoria}} \\
\cmidrule(lr){2-3}\cmidrule(lr){4-5}
\textbf{Medida} & \textbf{AMAWeb} & \textbf{AccessMonitor} & \textbf{AMAWeb} & \textbf{AccessMonitor} \\
\midrule
$n$                          &  116 &  123 &  105 &  122 \\
Média                        & 7,65 & 7,01 & 8,00 & 7,49 \\
Mediana                      & 7,70 & 7,20 & 7,90 & 7,30 \\
Desvio-padrão                & 1,00 & 1,22 & 0,70 & 0,91 \\
Mínimo                       & 4,20 & 3,80 & 6,30 & 4,80 \\
Máximo                       & 9,80 & 9,50 & 9,90 & 9,50 \\
\% acima do ponto de corte   &   XX &   XX &   XX &   XX \\
\bottomrule
\end{tabular}
\end{table}

%talvez fazer um gráfico boxplot

Nas duas ferramentas, as ouvidorias obtiveram média ligeiramente superior à dos portais e menor dispersão. Considerando apenas os municípios com os dois ambientes avaliados, a ouvidoria obteve nota maior que o portal em 55 dos 98 municípios no AMAWeb, com diferença média de $-0{,}41$ ponto. No AccessMonitor, isso ocorreu em 55 dos 114 municípios, com diferença média de $-0{,}51$ ponto. A diferença absoluta entre as notas do portal e da ouvidoria foi, em média, de 0,92 e 1,19 ponto, respectivamente. Assim, a nota do portal não substitui a avaliação específica da ouvidoria.

A comparação por tipo de solução, na Tabela~\ref{tab:tipos}, mostra diferenças entre os grupos. Nas ouvidorias integradas ao portal e nos sistemas próprios, a diferença média entre as notas dos dois ambientes é próxima de zero. Já as ouvidorias terceirizadas obtiveram notas mais altas que o respectivo portal em 20 de 22 municípios no AMAWeb e em 23 de 25 no AccessMonitor, com diferenças médias de 1,63 e 2,50 pontos. Os dados mostram uma associação entre o tipo de ouvidoria e a diferença de notas, mas não permitem atribuí-la ao fornecedor ou à política de acessibilidade do município. No caso do Fala.BR, parte dos municípios utiliza a mesma página de registro, como observado na caracterização da coleta; por isso, suas notas não representam implementações independentes.

Nas 218 URLs de portais de prefeituras e de ouvidorias com nota válida nas duas ferramentas, o AMAWeb atribuiu, em média, 0,63 ponto a mais que o AccessMonitor. Essa diferença de nível justifica reportar as ferramentas separadamente.

Os erros mais frequentes, apurados no AMAWeb, estão listados na Tabela~\ref{tab:erros_portais}.

\begin{table}[H]
\centering
\caption{Práticas classificadas como não aceitáveis com maior frequência no AMAWeb}
\label{tab:erros_portais}
\small
\begin{tabular}{lccrr}
\toprule
 & & & \multicolumn{2}{c}{\textbf{Páginas com o erro (\%)}} \\
\cmidrule(lr){4-5}
\textbf{Erro} & \textbf{Critério} & \textbf{Nível} & \textbf{Portais} & \textbf{Ouvidorias} \\
 & \textbf{WCAG} & & ($n=116$) & ($n=105$) \\
\midrule
Contraste de cor insuficiente             & 1.4.3  & AA  & 84,5 & 73,3 \\
Links sem nome acessível                  & 4.1.2  & A   & 75,9 & 58,1 \\
Atributos \texttt{id} repetidos           & 4.1.1  & A   & 65,5 & 18,1 \\
Links de imagem sem texto alternativo     & 2.4.4  & A   & 58,6 & 54,3 \\
Ausência de \texttt{h1} ou \texttt{h1} duplicado & 1.3.1 & A & 50,9 & 57,1 \\
Sem link para saltar ao conteúdo          & 2.4.1  & A   & 45,7 & 75,2 \\
Foco removido por JavaScript              & 2.4.7  & AA  & 41,4 & 71,4 \\
Hierarquia de cabeçalhos violada          & 2.4.10 & AAA & 40,5 & 62,9 \\
Idioma principal marcado incorretamente   & 3.1.1  & A   &  5,2 & 53,3 \\
\bottomrule
\end{tabular}
\end{table}

O contraste de cor insuficiente é o problema mais comum nos dois ambientes, presente em 84,5\% dos portais e 73,3\% das ouvidorias, seguido por links sem nome acessível. Os perfis de erro, contudo, diferem. Nos portais predominam problemas de marcação, como \texttt{id} repetidos (65,5\% contra 18,1\%). Nas ouvidorias predominam problemas de navegação e de foco, como a ausência de mecanismo para saltar ao conteúdo principal (75,2\%) e a remoção do foco por JavaScript (71,4\%), barreiras que afetam diretamente quem navega por teclado ou leitor de tela ao preencher um formulário. O idioma principal marcado incorretamente aparece em 53,3\% das ouvidorias e em apenas 5,2\% dos portais, o que é compatível com a replicação de uma mesma plataforma em diversos municípios. O AccessMonitor reproduziu o mesmo padrão geral: contraste em 74,8\% dos portais e 68,9\% das ouvidorias, e ausência de link para saltar ao conteúdo em 69,7\% das ouvidorias.

\subsection{Relação com porte populacional e indicadores socioeconômicos}

A Tabela~\ref{tab:correlacao} apresenta os coeficientes de correlação de
Spearman entre as notas e os indicadores municipais.
 
\begin{table}[H]
\centering
\caption{Correlação de Spearman ($\rho$) entre notas e indicadores municipais}
\label{tab:correlacao}
\begin{tabular}{lrr}
\toprule
\textbf{Indicador} & \textbf{Portal} & \textbf{Ouvidoria} \\
\midrule
População (IBGE, 2026)   & XX & XX \\
PIB \textit{per capita}  & XX & XX \\
IDHM (2010)              & XX & XX \\
\bottomrule
\multicolumn{3}{l}{\footnotesize * $p < 0{,}05$; ** $p < 0{,}01$}
\end{tabular}
\end{table}
 
\begin{table}[ht]
\centering
\caption{Notas médias por faixa populacional}
\label{tab:faixas}
\begin{tabular}{lrrr}
\toprule
\textbf{Faixa populacional} & \textbf{$n$} & \textbf{Portal} & \textbf{Ouvidoria} \\
\midrule
Até 10 mil habitantes      & XX & XX & XX \\
De 10 a 50 mil habitantes  & XX & XX & XX \\
Acima de 50 mil habitantes & XX & XX & XX \\
\bottomrule
\end{tabular}
\end{table}
 


\subsection{Ranqueamento dos municípios}

%por ferramenta ou fazer a média entre as duas (Trazer os 10+ e 10-)

O ranqueamento considerou a média simples entre as notas do portal e da
ouvidoria. A Tabela~\ref{tab:ranking} apresenta os dez municípios mais bem
posicionados; a lista completa está disponível em XX. % TODO: Adicionar tabela completa no zenodo
 
\begin{table}[ht]
\centering
\caption{Dez municípios com maiores médias de acessibilidade}
\label{tab:ranking}
\begin{tabular}{rlrrrr}
\toprule
\textbf{\#} & \textbf{Município} & \textbf{Portal} & \textbf{Ouvidoria} & \textbf{Média} & \textbf{Diferença} \\
\midrule
1  & XX & XX & XX & XX & XX \\
2  & XX & XX & XX & XX & XX \\
3  & XX & XX & XX & XX & XX \\
4  & XX & XX & XX & XX & XX \\
5  & XX & XX & XX & XX & XX \\
6  & XX & XX & XX & XX & XX \\
7  & XX & XX & XX & XX & XX \\
8  & XX & XX & XX & XX & XX \\
9  & XX & XX & XX & XX & XX \\
10 & XX & XX & XX & XX & XX \\
\bottomrule
\end{tabular}
\end{table}



\begin{table}[H]
\centering
\caption{Dez municípios com menores médias de acessibilidade}
\label{tab:ranking_ultimos}
\begin{tabular}{rlrrrr}
\toprule
\textbf{\#} & \textbf{Município} & \textbf{Portal} & \textbf{Ouvidoria} & \textbf{Média} & \textbf{Diferença} \\
\midrule
143  & XX & XX & XX & XX & XX \\
142  & XX & XX & XX & XX & XX \\
141  & XX & XX & XX & XX & XX \\
140  & XX & XX & XX & XX & XX \\
139  & XX & XX & XX & XX & XX \\
138  & XX & XX & XX & XX & XX \\
137  & XX & XX & XX & XX & XX \\
136  & XX & XX & XX & XX & XX \\
135  & XX & XX & XX & XX & XX \\
134 & XX & XX & XX & XX & XX \\
\bottomrule
\end{tabular}
\end{table}

\section{Considerações finais}

\bibliographystyle{sbc}
\bibliography{sbc-template}

\end{document}
